\section{Reflection on Instructional Choices}\label{sec:id-reflection}
\owner{Group}
\pages{1}

This section records the choices the group is least sure of, and what would change them. The reflection is written before the evaluation, so it states the expectations the group has now and will check against the results.

\subsection{The Pass Mark for the Gate}
The four-out-of-five pass mark is a judgement. It may be too lenient, if a student can pass the gate while still holding M2, or too strict, if a student who understands the rule is routed to a lesson because of one slip. The usability study will show which of these happens, and the item set will be revised if needed.

\subsection{Hints and Productive Struggle}
The hints reduce the difficulty of each step, and that is useful for a novice. There is a risk that a student who receives the third hint learns the answer without doing the reasoning. The mastery standard counts hints against a student's record for this reason. Whether the standard is strict enough is a question for the evaluation.

\subsection{Learning Styles}
The app does not adapt to learning styles. This was a deliberate choice, based on the evidence in Section~\ref{sec:id-multimodal}. The group recognises that some students will prefer one representation, and the app's use of three representations for each idea is intended to serve them without sorting them.

\subsection{Relatedness and Motivation}
The app does not support relatedness, which is one of the basic needs in self-determination theory. The app is used alone, and the group did not add a social feature because it would have added a shared record that is outside the scope of the prototype. This is the clearest gap in the motivational design.

\subsection{Scope}
The group chose to leave out deletion and self-balancing trees. The choice keeps the first exposure to trees focused on one rule, but it means the app does not teach the two-child case that is in the lecture. The choice is recorded in the SE document as decision D3, and the group would revisit it if the lecture plan changed.

\subsection{What the Group Would Change}
Two things would change first in a second version. The first is a larger item bank with items generated from the tree structure, so that the diagnostic cannot be memorised. The second is a check on whether the lesson order matches the lecture order, since the app's sequence was chosen by the group and not by the lecturer. \todo{add the group's findings from the usability study once it is run}
